\import{set/equinumerosity.tex}

\section{Equinumerosity and disjoint unions}

\begin{theorem}\label{disjoint_union_equinum_congruent}
    Suppose $A\approx C$ and $B\approx D$.
    Suppose $A\inter B=\emptyset$ and $C\inter D=\emptyset$.
    Then $A\union B\approx C\union D$.
\end{theorem}
\begin{proof}
    Take $f$ such that $f\in\bijections{A}{C}$ by \obv{equinum}.
    Take $g$ such that $g\in\bijections{B}{D}$ by \obv{equinum}.
    $f\in\funs{A}{C}$ and $g\in\funs{B}{D}$
        by \obv{bijections_to_funs}.
    Hence $f$ and $g$ are functions and $\dom{f}=A$ and $\dom{g}=B$
        by \obv{funs_is_function,funs}.
    Hence $f\in\funs{A}{C\union D}$ and
        $g\in\funs{B}{C\union D}$
        by \obv{funs_weaken_codom,union_upper_left,union_upper_right}.
    $f$ and $g$ are injective by \obv{bijections}.
    Hence $f\in\inj{A}{C\union D}$ and
        $g\in\inj{B}{C\union D}$
        by \auto{inj,funs_is_function,injective_function}.
    $\ran{f}=C$ and $\ran{g}=D$
        by \auto{bijections,ran_of_surj}.
    \begin{align*}
        \ran{f}\inter\ran{g}
            &= C\inter D
            \\
            &= \emptyset
    \end{align*}
    Hence $f\union g\in\inj{A\union B}{C\union D}$
        by \obv{union_inj_disjoint}.
    Hence $f\union g\in\funs{A\union B}{C\union D}$
        by \obv{inj}.
    Hence $f\union g$ is a function and
        $\dom{f\union g}=A\union B$
        by \obv{funs_is_function,funs}.
    $\ran{f\union g}=\ran{f}\union\ran{g}$ by \obv{ran_union}.
    Hence $\ran{f\union g}=C\union D$.
    Hence $f\union g\in\surj{A\union B}{C\union D}$
        by \obv{funs_surj_iff}.
    $f\union g$ is injective
        by \auto{inj,injective_function}.
    Hence $f\union g\in\bijections{A\union B}{C\union D}$
        by \obv{bijections}.
    Follows by \obv{equinum}.
\end{proof}

\begin{theorem}\label{three_disjoint_unions_equinum_congruent}
    Suppose $A\approx D$ and $B\approx E$ and $C\approx F$.
    Suppose $A\inter B=\emptyset$ and $A\inter C=\emptyset$ and
        $B\inter C=\emptyset$.
    Suppose $D\inter E=\emptyset$ and $D\inter F=\emptyset$ and
        $E\inter F=\emptyset$.
    Then $(A\union B)\union C\approx(D\union E)\union F$.
\end{theorem}
\begin{proof}
    $A\union B\approx D\union E$
        by \obv{disjoint_union_equinum_congruent}.
    $(A\union B)\inter C
        =(A\inter C)\union(B\inter C)$
        by \auto{inter_comm,inter_distrib_union}.
    Hence $(A\union B)\inter C=\emptyset$
        by \auto{union_emptyset}.
    $(D\union E)\inter F
        =(D\inter F)\union(E\inter F)$
        by \auto{inter_comm,inter_distrib_union}.
    Hence $(D\union E)\inter F=\emptyset$
        by \obv{union_emptyset}.
    Follows by \obv{disjoint_union_equinum_congruent}.
\end{proof}
